% Propuesta: acondicionamiento de señal para micrófonos IEPE (sonómetro de un canal)
% Compilar con:  tectonic propuesta_sonometro.tex
\documentclass[11pt,a4paper]{article}

\usepackage{fontspec}
\usepackage[spanish,es-noshorthands,es-tabla]{babel}
\usepackage[margin=2.3cm]{geometry}
\usepackage{amsmath}
\usepackage{siunitx}
\sisetup{output-decimal-marker={,}, group-separator={.}, per-mode=symbol}
\usepackage{booktabs,tabularx,array,multirow}
\usepackage{graphicx}
\usepackage{xcolor}
\usepackage{enumitem}
\setlist{itemsep=2pt,topsep=4pt}
\usepackage{float}
\usepackage{caption}
\usepackage{tikz}
\usetikzlibrary{positioning,arrows.meta,calc,babel}
\usepackage[american,siunitx]{circuitikz}
\usepackage{pgfplots}
\pgfplotsset{compat=1.17}
\usepackage{listings}
\lstset{basicstyle=\ttfamily\small, frame=single, columns=fullflexible, keepspaces=true, showstringspaces=false}
\usepackage{tcolorbox}
\usepackage{xurl}
\usepackage[hidelinks]{hyperref}

\newcolumntype{Y}{>{\raggedright\arraybackslash}X}
\newcommand{\nota}[1]{\textcolor{red!60!black}{\textbf{[#1]}}}

\title{\textbf{Propuesta de acondicionamiento de señal\\para dos micrófonos IEPE}\\[4pt]
\large Sonómetro de un canal con selección de micrófono}
\author{Grupo: \nota{nombres} \\ Curso: \nota{nombre del curso} --- Prof.\ \nota{nombre}}
\date{4 de noviembre de 2026}

\begin{document}
\maketitle

\begin{tcolorbox}[title=Resumen de la propuesta]
Se propone alimentar los dos micrófonos disponibles (preamplificador \textbf{SVANTEK SV~12} y micrófono \textbf{ROGA RG-50}) con
\textbf{fuentes de corriente constante} independientes construidas con componentes de paso (TL431 + transistor PNP),
desde una fuente de \SI{24}{\volt}. Ambos micrófonos son de tipo \textbf{IEPE}: la información acústica viaja como una
\textbf{tensión alterna montada sobre una tensión continua de polarización} ($\approx$\SIrange{8}{14}{\volt}) en el mismo
cable coaxial que lleva la corriente de alimentación. La señal se extrae con un \textbf{condensador de acople}, se lleva a un
nivel centrado en \SI{1.65}{\volt}, se amplifica ($G=11$), se filtra y se observa en (1) el \textbf{osciloscopio} del
laboratorio y (2) el \textbf{ADC de un microcontrolador} que envía la forma de onda a un PC. Un conmutador permite
elegir el micrófono. El montaje se hace sin PCB, con componentes conseguibles en Manizales (Electrónicos Caldas) y,
para lo que falte, en Medellín (Didácticas Electrónicas). Costo estimado de componentes: \textbf{$\approx$ \$40.000--\$50.000~COP}
(sin contar fuente de laboratorio ni osciloscopio).
\end{tcolorbox}

\tableofcontents
\newpage

%=====================================================================
\section{Contexto, objetivo y restricciones}

\textbf{Objetivo.} Diseñar y montar la etapa de entrada de un sonómetro de un canal que permita:
\begin{enumerate}
  \item alimentar correctamente dos micrófonos de medición distintos y \textbf{seleccionar} cuál se usa;
  \item \textbf{extraer} la señal acústica del cable de alimentación;
  \item \textbf{visualizar} la forma de onda en el osciloscopio y en un microcontrolador.
\end{enumerate}
El cálculo de nivel de presión sonora (dB SPL, ponderación A, $L_{eq}$) queda como extensión opcional (sección~\ref{sec:extension}).

\textbf{Restricciones.}
\begin{itemize}
  \item Fecha de entrega: \textbf{25 de noviembre de 2026} ($\approx$ 3 semanas).
  \item No se puede fabricar PCB. Montaje en protoboard y/o sobre placa de cobre virgen (técnica ``Manhattan'').
  \item Componentes comprados localmente (Manizales); preferencia por encapsulados de paso (DIP, TO-92).
  \item Instrumentación disponible: osciloscopio de laboratorio, fuente de laboratorio, microcontrolador.
\end{itemize}

%=====================================================================
\section{Los transductores}

\subsection{Identificación}

\textbf{SVANTEK SV 12 (n.º de serie 18210, recuperado).} En la nomenclatura de Svantek, ``SV~xx'' designa accesorios
y ``SVAN~9xx'' los equipos. El SV~12 / SV~12L es el \textbf{preamplificador IEPE de media pulgada} que usan los sonómetros
SVAN~94x, SVAN~958/958A y SVAN~977~\cite{svan977w,svan958a,sv977d}. Es decir, \emph{no es un micrófono completo}: la cápsula
(prepolarizada, p.\,ej.\ MK~255 o ACO~7052E) se enrosca encima. El número 18210 es el serial del preamplificador.
Se identifica por un conector \textbf{TNC} macho con rosca en la base. Si el conector fuera LEMO de 7 pines se trataría
de un SV~17 o SV~18, que \emph{no} son IEPE~\cite{sv971a}. \nota{Verificar: tipo de conector y si tiene cápsula montada.}

\textbf{ROGA Instruments RG-50.} Micrófono de $1/4''$ de tipo electret prepolarizado con electrónica IEPE integrada,
clase~1, salida por conector \textbf{BNC} hembra~\cite{rg50}. ROGA fabrica modelos casi idénticos (MI17, MI19); la
diferencia está en el grabado del cuerpo~\cite{roga14}.

\subsection{Especificaciones relevantes}

\begin{table}[H]
\centering
\caption{Especificaciones de alimentación y salida (fuentes: manuales Svantek y hoja de datos ROGA).}
\label{tab:mics}
\begin{tabularx}{\textwidth}{lYY}
\toprule
Parámetro & SVANTEK SV 12 / SV 12L & ROGA RG-50 \\
\midrule
Tipo & Preamplificador IEPE $1/2''$ para cápsulas prepolarizadas & Micrófono electret $1/4''$ con IEPE integrado \\
Corriente de alimentación & \SI{1.5}{\milli\ampere} (SVAN~977) y \SI{2.5}{\milli\ampere} (SVAN~958A); máx.\ no publicado & \SIrange{2}{20}{\milli\ampere} (versiones antiguas: \SIrange{2}{6}{\milli\ampere}) \\
Tensión de alimentación & hasta \SI{28}{\volt}--\SI{30}{\volt} & \SIrange{18}{30}{\volt} \\
Tensión DC de polarización & \textbf{No publicada} (se medirá) & \SI{12.5}{\volt} \\
Sensibilidad & La de la cápsula (\SIrange{20}{50}{\milli\volt\per\pascal}) $-$\SI{0.7}{\decibel} & \SI{50}{\milli\volt\per\pascal} nominal \\
Respuesta en frecuencia & Depende de la cápsula & \SI{5}{\hertz}--\SI{30}{\kilo\hertz} \\
Nivel máximo & $\geq$\SI{10.8}{\volt} pico a la salida & \SI{130}{\decibel} pico ($\approx$\SI{3.2}{\volt} pico) \\
Conector & TNC (centro: señal + alimentación; malla: tierra) & BNC hembra \\
Polarización 200 V / phantom 48 V & \textbf{No} (nunca aplicar) & \textbf{No} (nunca aplicar) \\
\bottomrule
\end{tabularx}
\end{table}

\textbf{Conclusión para el diseño:} una corriente de \SI{2.5}{\milli\ampere} está dentro de lo documentado por ambos
fabricantes. Aun así, como el grupo planteó usar corrientes distintas, la propuesta incluye una fuente de
\SI{2.5}{\milli\ampere} para el SV~12 y otra de \SI{4}{\milli\ampere} (valor típico de la industria) para el RG-50.

%=====================================================================
\section{¿Cómo viaja la señal en un micrófono IEPE?}
\label{sec:iepe}

IEPE (\emph{Integrated Electronics Piezo-Electric}; también llamado ICP, CCLD o DeltaTron) es un esquema de \textbf{dos hilos}:
un solo cable coaxial lleva la alimentación \emph{y} la señal~\cite{pcbguide,hbk}.

\begin{enumerate}
  \item El acondicionador entrega una \textbf{corriente constante} $I_0$ (típicamente \SIrange{2}{20}{\milli\ampere}) desde
        una fuente de \SIrange{18}{30}{\volt}.
  \item La electrónica interna del micrófono fija en el conector una \textbf{tensión DC de polarización} $V_B$ (\SIrange{8}{14}{\volt}).
  \item La presión acústica $p(t)$ produce una variación de tensión que se \textbf{suma} a $V_B$:
\end{enumerate}
\begin{equation}
  v_{\text{nodo}}(t) = V_B + S\cdot p(t),
  \label{eq:iepe}
\end{equation}
donde $S$ es la sensibilidad en \si{\milli\volt\per\pascal}. \textbf{La información viaja como tensión, no como corriente}:
la corriente es constante y solo sirve para alimentar. Para recuperar la señal basta eliminar la componente DC con un
condensador (acople AC). La figura~\ref{fig:senal} lo ilustra.

\begin{figure}[H]
\centering
\begin{tikzpicture}
\begin{axis}[name=a1, width=7.2cm, height=4.6cm, xlabel={$t$ [ms]}, ylabel={$v_{\text{nodo}}$ [V]},
  ymin=0, ymax=16, domain=0:3, samples=300, title={En el conector BNC}, grid=major]
  \addplot[blue, thick] {12.5 + 1.0*sin(deg(2*pi*x))};
  \addplot[red, dashed] {12.5};
  \node[anchor=south west, font=\small] at (axis cs:0.05,13.6) {$V_B$};
\end{axis}
\begin{axis}[at={(a1.east)}, anchor=west, xshift=1.6cm, width=7.2cm, height=4.6cm, xlabel={$t$ [ms]},
  ylabel={$v$ [V]}, ymin=-2, ymax=2, domain=0:3, samples=300, title={Después del condensador}, grid=major]
  \addplot[blue, thick] {1.0*sin(deg(2*pi*x))};
\end{axis}
\end{tikzpicture}
\caption{Señal IEPE (tono de \SI{1}{\kilo\hertz}, amplitud exagerada): DC de polarización + señal acústica, y la misma señal tras el acople AC.}
\label{fig:senal}
\end{figure}

\textbf{Órdenes de magnitud} con $S=\SI{50}{\milli\volt\per\pascal}$ (RG-50):
\begin{itemize}
  \item Voz a \SI{1}{\metre} ($\approx$\SI{60}{\decibel} SPL, \SI{0.02}{\pascal}): $\approx$\SI{1}{\milli\volt} rms.
  \item Parlante de celular cerca ($\approx$\SIrange{80}{90}{\decibel}): $\approx$\SIrange{5}{16}{\milli\volt} rms.
  \item Calibrador acústico \SI{94}{\decibel} (\SI{1}{\pascal}): \SI{50}{\milli\volt} rms.
\end{itemize}

\textbf{Longitud de cable.} PCB Piezotronics da la frecuencia máxima que la corriente puede excitar en la capacidad del cable~\cite{pcbcable}:
\begin{equation}
  f_{\max} = \frac{10^9}{2\pi\, C\, \dfrac{V_p}{I_c - 1}}\quad [\si{\hertz}],\qquad C\ [\si{\pico\farad}],\ V_p\ [\si{\volt}],\ I_c\ [\si{\milli\ampere}].
\end{equation}
Con \SI{10}{\metre} de RG-58 ($\approx$\SI{1000}{\pico\farad}), $V_p=\SI{3.2}{\volt}$ e $I_c=\SI{2.5}{\milli\ampere}$:
$f_{\max}\approx\SI{75}{\kilo\hertz}$, muy por encima de \SI{20}{\kilo\hertz}. Con cables de banco ($\leq$\SI{2}{\metre}) no hay
ninguna limitación. \textbf{La corriente no necesita ser grande ni precisa}; TI lo resume así: ``la exactitud de la
fuente de corriente no es crítica, pero sí el rango de compliancia, el ruido y la regulación''~\cite{tida010249}.

%=====================================================================
\section{Arquitectura propuesta}

\begin{figure}[H]
\centering
\begin{tikzpicture}[
  blk/.style={draw, rounded corners, thick, align=center, minimum height=1.05cm, minimum width=2.2cm, font=\small},
  >={Stealth}]
  \node[blk, fill=gray!10] (ps) at (0,0) {Fuente\\\SI{24}{\volt}};
  \node[blk, fill=blue!8] (ia) at (2.9,1.1) {Fuente I\\\SI{2.5}{\milli\ampere}};
  \node[blk, fill=blue!8] (ib) at (2.9,-1.1) {Fuente I\\\SI{3.9}{\milli\ampere}};
  \node[blk, fill=green!8] (ma) at (5.8,1.1) {SV 12\\(TNC$\to$BNC)};
  \node[blk, fill=green!8] (mb) at (5.8,-1.1) {RG-50\\(BNC)};
  \node[blk, fill=orange!10] (sel) at (8.7,0) {Acople AC\\+ selector};
  \node[blk, fill=orange!10] (amp) at (11.6,0) {Ganancia\\$G=11$\\+ $V_{MID}$};
  \node[blk, fill=orange!10] (aa) at (11.6,-2.9) {Filtro\\antialias};
  \node[blk, fill=purple!8] (adc) at (8.7,-2.9) {ADC\\(Pico / ESP32)};
  \node[blk, fill=purple!8] (pc) at (5.8,-2.9) {PC\\(Python)};
  \node[blk, fill=red!8] (osc) at (11.6,2.6) {Osciloscopio};
  \draw[->] (ps) -- (ia); \draw[->] (ps) -- (ib);
  \draw[->] (ia) -- (ma); \draw[->] (ib) -- (mb);
  \draw[->] (ma) -- (sel); \draw[->] (mb) -- (sel);
  \draw[->] (sel) -- (amp); \draw[->] (amp) -- (aa); \draw[->] (aa) -- (adc); \draw[->] (adc) -- (pc);
  \draw[->, dashed] (ma.north) |- (osc.west) node[pos=0.75, above, font=\scriptsize]{BNC directo, acople AC};
  \draw[->, dashed] (amp) -- (osc);
\end{tikzpicture}
\caption{Diagrama de bloques. Las líneas punteadas son los puntos de observación con el osciloscopio.}
\label{fig:bloques}
\end{figure}

El desarrollo se plantea en dos hitos que se pueden demostrar por separado:
\begin{description}
  \item[Hito 1 (mínimo viable):] fuente de corriente + micrófono + osciloscopio con acople AC. Basta para ``ver la forma de
        onda'' y no requiere amplificador.
  \item[Hito 2:] etapa de acople, ganancia y filtro hacia el ADC del microcontrolador, y gráfica en el PC.
\end{description}

%=====================================================================
\section{Fuente de corriente constante}

\subsection{Opciones evaluadas}

\begin{table}[H]
\centering
\small
\caption{Comparación de fuentes de corriente para IEPE (\SIrange{2}{20}{\milli\ampere}, compliancia \SIrange{18}{30}{\volt}).}
\label{tab:fuentes}
\begin{tabularx}{\textwidth}{>{\raggedright\arraybackslash}p{2.6cm}YY>{\raggedright\arraybackslash}p{2.6cm}}
\toprule
Opción & Ventajas & Desventajas & Disponibilidad \\
\midrule
\textbf{TL431 + PNP} (elegida) & Pocos componentes de paso, caída $\approx$\SI{2}{\volt}, corriente fija por una sola resistencia; funciona igual a \SI{2.5}{} y \SI{4}{\milli\ampere} & Exactitud $\pm$\SIrange{3}{5}{\percent} (suficiente para IEPE) & TL431 en stock local; PNP 2N3906/BC557 comunes \\
LM334 & El clásico; $I=\SI{67.7}{\milli\volt}/R$, caída $\approx$\SI{1}{\volt}~\cite{lm334} & Deriva $+$\SI{0.34}{\percent\per\celsius}; máx.\ \SI{10}{\milli\ampere} & \textbf{No disponible localmente} (pedido especial 10--16 días) \\
LM317 como fuente & Muy común; $I=\SI{1.25}{\volt}/R$; usado por Endevco~\cite{endevco} & Requiere corriente mínima de carga \SIrange{3.5}{10}{\milli\ampere}~\cite{lm317}: \textbf{no sirve para \SI{2.5}{\milli\ampere}} & En stock local \\
JFET + resistencia & Muy simple & $I_{DSS}$ varía 3--5$\times$ entre unidades: hay que ajustar cada una & J111/J112 en stock \\
Diodo regulador (CRD, E-352) & Lo que usa PCB en sus acondicionadores~\cite{pcbguide} & Valor fijo; difícil de conseguir & No disponible \\
LT3092 (ADI) & Referencia de ADI en CN0540/CN0582~\cite{cn0540,cn0582} & Solo SMD & No disponible \\
Op-amp + transistor & Máxima exactitud (TI TIDA-010249)~\cite{tida010249,slyt768} & Más complejo; requiere op-amp \emph{rail-to-rail} & Parcial \\
\bottomrule
\end{tabularx}
\end{table}

\subsection{Circuito elegido: referencia TL431 + transistor PNP}

El TL431 conectado como referencia (pin REF unido al cátodo) mantiene \SI{2.495}{\volt} entre $V_{CC}$ y la base del PNP.
El emisor queda una caída $V_{EB}\approx\SI{0.65}{\volt}$ por encima de la base, de modo que sobre $R_E$ hay una tensión
constante y, por tanto, una corriente constante que sale por el colector hacia el micrófono (figura~\ref{fig:fuente}).
Es la misma idea de ``referencia + transistor'' que analiza TI en~\cite{slyt768}.

\begin{figure}[tbp]
\centering
\begin{circuitikz}[scale=0.95, transform shape]
  \coordinate (top) at (0,7);
  \draw (-0.5,7) node[left]{$V_{CC}=\SI{24}{\volt}$} -- (7.5,7);
  % TL431 (cátodo arriba)
  \draw (2,4) to[zD, l={TL431 (REF unido a K)}, *-*] (2,7);
  \draw (2,4) to[R, l_={$R_Z=\SI{10}{\kilo\ohm}$}] (2,1) node[ground]{};
  % PNP
  \node[pnp, anchor=B] (Q) at (4.2,4) {};
  \node[right=0.05cm of Q.text, font=\small] {2N3906 / BC557};
  \draw (2,4) -- (Q.B);
  \draw (Q.E |- top) to[R, l={$R_E$}, *-] (Q.E);
  \draw (Q.C) -- ++(0,-1.0) coordinate (N);
  \draw (N) to[short, *-o, i>^={$I_0$}] ++(3,0) node[right, align=left]{Nodo IEPE\\(centro BNC)};
  \draw (N) -- ++(0,-0.4) node[below, font=\small]{\textit{a $C_1$ (fig.~\ref{fig:acople})}};
  % desacople
  \draw (7.2,7) to[C, l={\SI{100}{\nano\farad}}, *-] (7.2,5.2) node[ground]{};
\end{circuitikz}
\caption{Fuente de corriente de un canal. Se construyen dos: $R_E=\SI{750}{\ohm}$ (\SI{2.5}{\milli\ampere}, SV 12) y
$R_E=\SI{470}{\ohm}$ (\SI{3.9}{\milli\ampere}, RG-50). La malla del BNC va a tierra.}
\label{fig:fuente}
\end{figure}

\textbf{Cálculos.}
\begin{align}
  I_0 &\approx \frac{V_{REF}-V_{EB}}{R_E} = \frac{\SI{2.495}{\volt}-\SI{0.65}{\volt}}{R_E} = \frac{\SI{1.845}{\volt}}{R_E} \\
  R_E(\SI{2.5}{\milli\ampere}) &= \SI{738}{\ohm}\ \Rightarrow\ \SI{750}{\ohm}\ (I_0=\SI{2.46}{\milli\ampere}) \\
  R_E(\SI{4}{\milli\ampere}) &= \SI{461}{\ohm}\ \Rightarrow\ \SI{470}{\ohm}\ (I_0=\SI{3.93}{\milli\ampere}) \\
  I_Z &= \frac{V_{CC}-V_{REF}}{R_Z} = \frac{\SI{21.5}{\volt}}{\SI{10}{\kilo\ohm}} = \SI{2.15}{\milli\ampere}\ \ (\geq \SI{1}{\milli\ampere}\text{ requerido por el TL431~\cite{tl431}})
\end{align}

\textbf{Compliancia y margen de señal.} La tensión máxima que puede alcanzar el nodo es
\begin{equation}
  V_{\text{nodo,máx}} \approx V_{CC} - (V_{REF}-V_{EB}) - V_{CE,sat} \approx 24 - 1{,}85 - 0{,}3 \approx \SI{21.9}{\volt}.
\end{equation}
Con $V_B=\SI{12.5}{\volt}$ (RG-50) quedan $\approx$\SI{9}{\volt} de excursión positiva, frente a \SI{3.2}{\volt} pico que
entrega el micrófono a \SI{130}{\decibel}. La disipación del PNP es $(V_E-V_B)\,I_0\approx 9{,}6\times 4\approx\SI{40}{\milli\watt}$,
muy por debajo de los \SI{625}{\milli\watt} del encapsulado TO-92.

\textbf{Observaciones.}
\begin{itemize}
  \item La exactitud ($\pm$\SIrange{3}{5}{\percent}, dominada por $V_{EB}$) es suficiente: los fabricantes de IEPE aceptan
        \SIrange{2}{20}{\milli\ampere}. Se mide la corriente real y, si hace falta, se ajusta $R_E$.
  \item No se coloca condensador directamente sobre el TL431 (ciertos valores lo vuelven inestable~\cite{tl431}).
  \item Si el micrófono se desconecta, el nodo sube a $\approx$\SI{22}{\volt} (circuito abierto) y la corriente cae a cero:
        la fuente \emph{falla de forma segura}. Nunca se conecta una tensión al micrófono sin limitación de corriente~\cite{pcbguide}.
  \item Alternativa si se consigue un LM334: $R=\SI{68}{\milli\volt}/I_0$ ($\SI{27}{\ohm}\to\SI{2.5}{\milli\ampere}$; $\SI{18}{\ohm}\to\SI{3.8}{\milli\ampere}$).
\end{itemize}

%=====================================================================
\section{Selección entre los dos micrófonos}

\begin{table}[H]
\centering
\small
\caption{Dos formas de seleccionar el micrófono.}
\begin{tabularx}{\textwidth}{lYY}
\toprule
 & \textbf{A. Dos fuentes, selector después del acople (recomendada)} & B. Una fuente con conmutador DPDT \\
\midrule
Funcionamiento & Cada micrófono tiene su fuente, su $C_1$ y su $R_B$, ambos siempre alimentados. Un interruptor SPDT elige qué señal (ya sin DC) entra al amplificador. & Un polo conmuta $R_E$ (\SI{750}{\ohm} / \SI{470}{\ohm}); el otro conmuta la salida de la fuente al BNC A o al BNC B. \\
Transitorio al conmutar & Ninguno: los condensadores ya están cargados & Salto de $\approx$\SI{12}{\volt} en $C_1$ y arranque del micrófono; hay que esperar unos segundos \\
Componentes extra & 1 TL431, 1 PNP, 3 resistencias, 1 condensador ($\approx$\$3.000 COP) & Conmutador DPDT \\
Consumo & $\approx$\SI{6.4}{\milli\ampere} + \SI{4.3}{\milli\ampere} ($R_Z$) & $\approx$\SI{4}{\milli\ampere} + \SI{2.2}{\milli\ampere} \\
\bottomrule
\end{tabularx}
\end{table}

\begin{figure}[H]
\centering
\begin{circuitikz}[scale=0.9, transform shape]
  \node[draw, rounded corners, thick, align=center, minimum width=2.4cm, minimum height=1cm] (FA) at (0,2) {Fuente A\\\SI{2.5}{\milli\ampere}};
  \node[draw, rounded corners, thick, align=center, minimum width=2.4cm, minimum height=1cm] (FB) at (0,-1) {Fuente B\\\SI{3.9}{\milli\ampere}};
  \draw (FA.east) -- (2,2) coordinate (NA);
  \draw (FB.east) -- (2,-1) coordinate (NB);
  \draw (NA) to[short, *-o] (2,3) node[above]{SV 12};
  \draw (NB) to[short, *-o] (2,0) node[above]{RG-50};
  \draw (NA) to[C, l={$C_{1A}$}] (4.5,2) coordinate (XA);
  \draw (NB) to[C, l={$C_{1B}$}] (4.5,-1) coordinate (XB);
  \draw (XA) to[R, l={$R_{BA}$}, *-] (4.5,0.4) node[below, font=\small]{$V_{MID}$};
  \draw (XB) to[R, l={$R_{BB}$}, *-] (4.5,-2.6) node[below, font=\small]{$V_{MID}$};
  \node[spdt, xscale=-1, anchor=in] (S) at (8,0.5) {};
  \draw (XA) -- (6.2,2) |- (S.out 1);
  \draw (XB) -- (6.2,-1) |- (S.out 2);
  \draw (S.in) -- ++(0.6,0) node[right, align=left]{a $R_S$ y\\amplificador};
\end{circuitikz}
\caption{Opción A: dos canales siempre alimentados y selector después del acople.}
\label{fig:selector}
\end{figure}

%=====================================================================
\section{Extracción y acondicionamiento de la señal}

\begin{figure}[H]
\centering
\begin{circuitikz}[scale=0.92, transform shape]
  % acople
  \draw (0,3) node[left]{Nodo IEPE} to[short, o-] (0.4,3) to[C, l={$C_1$}] (2.6,3) coordinate (A);
  \draw (A) to[R, l_={$R_B$}, *-] (2.6,1.1) node[below]{$V_{MID}$};
  \draw (A) to[R, l={$R_S$}] (5,3) coordinate (P);
  % protección
  \draw (P) to[D, *-] (5,4.6) node[above]{$+\SI{3.3}{\volt}$};
  \draw (5,1.4) node[ground]{} to[D] (P);
  % op-amp
  \node[op amp, noinv input up, anchor=+] (U) at (6.2,3) {};
  \draw (P) -- (U.+);
  \draw (U.-) -- ++(-0.4,0) -- ++(0,-2.1) coordinate (M);
  \draw (M) to[R, l_={$R_G$}] ++(-2,0) node[left]{$V_{MID}$};
  \draw (M) to[R, l={$R_F$}, *-] (M -| U.out) coordinate (Fo);
  \draw (Fo) -- (U.out);
  \draw (M) -- ++(0,-1.2) coordinate (M2) to[C, l_={$C_F$}] (M2 -| Fo) -- (Fo);
  % salida
  \draw (U.out) to[short, *-] ++(0.6,0) coordinate (O) to[R, l={$R_O$}] ++(2.2,0) coordinate (ADC) to[short, -o] ++(0.9,0) node[right]{ADC};
  \draw (ADC) to[C, l={$C_O$}, *-] ++(0,-1.8) node[ground]{};
  \draw (O) to[short, -o] ++(0,1.2) node[above]{Osciloscopio};
\end{circuitikz}
\caption{Acople AC, protección, ganancia no inversora y filtro antialias. Op-amp: MCP6002 alimentado a \SI{3.3}{\volt} (o LM358 a \SI{5}{\volt}).}
\label{fig:acople}
\end{figure}

\begin{table}[H]
\centering
\small
\caption{Valores propuestos para la figura~\ref{fig:acople}.}
\begin{tabularx}{\textwidth}{llY}
\toprule
Componente & Valor & Función \\
\midrule
$C_1$ & \SI{1}{\micro\farad} poliéster, $\geq$\SI{63}{\volt} & Bloquea la DC ($V_B$). Debe soportar $\approx$\SI{22}{\volt} en circuito abierto \\
$R_B$ & \SI{100}{\kilo\ohm} & Referencia la señal a $V_{MID}=\SI{1.65}{\volt}$ (en vez de tierra) \\
$R_S$ + 2$\times$1N4148 & \SI{4.7}{\kilo\ohm} & Limita la corriente cuando se conecta un micrófono en caliente (salto de $\pm$\SI{12}{\volt}) \\
$R_F$, $R_G$ & \SI{10}{\kilo\ohm}, \SI{1}{\kilo\ohm} & Ganancia $G=11$ \\
$C_F$ & \SI{470}{\pico\farad} & Polo en \SI{34}{\kilo\hertz} \\
$R_O$, $C_O$ & \SI{1}{\kilo\ohm}, \SI{4.7}{\nano\farad} & Antialias de 1.\textsuperscript{er} orden (\SI{34}{\kilo\hertz}) y protección del pin ADC \\
$V_{MID}$ & \SI{10}{\kilo\ohm}/\SI{10}{\kilo\ohm} + \SI{10}{\micro\farad} + seguidor & Tensión de media escala con baja impedancia (segunda mitad del op-amp) \\
\bottomrule
\end{tabularx}
\end{table}

\textbf{Cálculos.}
\begin{align}
  f_{c,\text{acople}} &= \frac{1}{2\pi R_B C_1} = \frac{1}{2\pi(\SI{100}{\kilo\ohm})(\SI{1}{\micro\farad})} = \SI{1.6}{\hertz}
  \qquad (\text{pérdida a } \SI{20}{\hertz}: \SI{0.03}{\decibel}) \\
  \tau &= R_B C_1 = \SI{0.1}{\second}\ \Rightarrow\ \text{estable tras } 5\tau \approx \SI{0.5}{\second} \\
  G &= 1 + \frac{R_F}{R_G} = 11, \qquad f_{p} = \frac{1}{2\pi R_F C_F} \approx \SI{34}{\kilo\hertz} \\
  v_{\text{ADC}}(t) &= V_{MID} + G\cdot S\cdot p(t)
\end{align}
Con el ADC de \SIrange{0}{3.3}{\volt} y $\pm$\SI{1.5}{\volt} útiles alrededor de \SI{1.65}{\volt}, la entrada máxima es
\SI{136}{\milli\volt} pico. Con \SI{50}{\milli\volt\per\pascal} el sistema satura hacia $\approx$\SI{100}{\decibel} SPL, y un
calibrador de \SI{94}{\decibel} produce \SI{0.78}{\volt} pico en el ADC. Para sonidos débiles se puede subir a $G=101$
($R_F=\SI{100}{\kilo\ohm}$).

\textbf{Elección del op-amp.} El TL072 y el NE5532 (en stock local) necesitan alimentación partida o mayor a
\SI{10}{\volt}; sirven para una salida al osciloscopio pero no para entregar \SIrange{0}{3.3}{\volt} al ADC. Se propone el
\textbf{MCP6002} (\emph{rail-to-rail}, \SI{1}{\mega\hertz}, DIP-8), alimentado desde los \SI{3.3}{\volt} del microcontrolador.
Alternativa local: \textbf{LM358} alimentado a \SI{5}{\volt}, cuya entrada y salida cubren \SIrange{0}{3.5}{\volt}.
Los diodos y $R_O$ impiden que el pin del ADC reciba más de \SI{3.3}{\volt}.

\textbf{Detección de fallas (opcional).} La tensión DC del nodo indica el estado del micrófono~\cite{hbk,tida010249}:
$<\SI{2}{\volt}$ = cortocircuito; \SIrange{7}{14}{\volt} = normal; $\approx$\SI{22}{\volt} = circuito abierto o micrófono dañado.
Basta un multímetro, o un divisor \SI{1}{\mega\ohm}/\SI{120}{\kilo\ohm} hacia otro canal del ADC.

%=====================================================================
\section{Visualización}

\subsection{Osciloscopio (Hito 1)}
Se conecta una T BNC en el nodo IEPE (o se mide en la salida del amplificador):
\begin{itemize}
  \item Acople \textbf{AC} (elimina los $\approx$\SI{12}{\volt} de DC), \SIrange{5}{20}{\milli\volt}/div en el nodo.
  \item Base de tiempo \SI{200}{\micro\second}/div para un tono de \SI{1}{\kilo\hertz}.
  \item Límite de ancho de banda (\SI{20}{\mega\hertz}) activado, disparo con rechazo de HF, promediado si hay ruido.
  \item Función \textbf{FFT} del osciloscopio para mostrar el pico de \SI{1}{\kilo\hertz} y el zumbido de \SI{60}{\hertz}.
\end{itemize}

\subsection{Microcontrolador (Hito 2)}

\begin{table}[H]
\centering
\small
\caption{Microcontroladores considerados.}
\begin{tabularx}{\textwidth}{>{\raggedright\arraybackslash}p{2.9cm}>{\raggedright\arraybackslash}p{2.8cm}Y>{\raggedright\arraybackslash}p{2.7cm}}
\toprule
Placa & ADC & Comentario & Disponibilidad \\
\midrule
\textbf{Raspberry Pi Pico} (recomendada) & 12 bits, \SI{500}{kS/s}, \SIrange{0}{3.3}{\volt} & Ejemplo oficial \texttt{adc/dma\_capture}~\cite{picoex}. ENOB $\approx$8,7 bits~\cite{rp2040err}: suficiente para ver la onda & Didácticas (\$19.900) \\
ESP32 DevKitC & 12 bits, modo continuo (DMA) & No lineal; ventana útil $\approx$\SIrange{0.15}{2.45}{\volt}: usar $V_{MID}\approx\SI{1.3}{\volt}$ y menor ganancia. Usar ADC1~\cite{espadc} & Caldas (\$27.300) \\
Arduino Uno & 10 bits, $\approx$\SI{10}{kS/s} & Solo ráfagas cortas a \SI{1}{\kilo\hertz}; 2~kB de RAM & Caldas \\
\bottomrule
\end{tabularx}
\end{table}

\textbf{Frecuencia de muestreo.} Para ver la forma de un tono de \SI{1}{\kilo\hertz} bastan 20--50 muestras por ciclo
(\SIrange{20}{50}{kS/s}). Para toda la banda de audio, Nyquist exige $f_s>\SI{40}{\kilo\hertz}$; con el
antialias de \SI{34}{\kilo\hertz} se recomienda sobremuestrear a \SIrange{100}{200}{kS/s} (posible en el Pico).

\textbf{Captura por ráfagas.} Un puerto serie a \SI{115200}{baud} transmite $\approx$\SI{11.5}{kB/s},
insuficiente para transmitir audio en tiempo real. Por eso el microcontrolador \textbf{captura un bloque} (p.\,ej.\ 4096
muestras) en RAM vía DMA y luego lo envía. En el PC, un script de Python grafica la forma de onda y su FFT:

\begin{lstlisting}[language=Python]
import serial, numpy as np, matplotlib.pyplot as plt
FS = 100_000                                   # frecuencia de muestreo [Hz]
s = serial.Serial("COM5", 115200, timeout=5)   # puerto del Pico
s.write(b"c")                                  # pedir una captura
x = np.array([int(l) for l in s.read_until(b"END").decode().split()[:-1]])
v = x * 3.3 / 4095                             # cuentas -> voltios
t = np.arange(len(v)) / FS * 1e3
f = np.fft.rfftfreq(len(v), 1/FS)
V = 20*np.log10(np.abs(np.fft.rfft((v - v.mean()) * np.hanning(len(v)))) + 1e-12)
fig, (a, b) = plt.subplots(2, 1)
a.plot(t, v); a.set_xlabel("t [ms]"); a.set_ylabel("V")
b.plot(f, V); b.set_xlabel("f [Hz]"); b.set_ylabel("dB")
plt.show()
\end{lstlisting}
Alternativas sin programar la interfaz: \emph{SerialPlot}~\cite{serialplot} o \emph{Scoppy} (osciloscopio con Pico + app Android)~\cite{scoppy}.

%=====================================================================
\section{Alimentación, tierra y ruido}

\begin{itemize}
  \item \textbf{Fuente:} fuente de laboratorio a \SI{24}{\volt} con límite de corriente en \SIrange{30}{50}{\milli\ampere}
        (o dos baterías de \SI{12}{\volt} en serie, la opción más silenciosa). Diodo 1N4007 en serie contra inversión de
        polaridad, filtro RC \SI{47}{\ohm} + \SI{100}{\micro\farad}/\SI{35}{\volt} + \SI{100}{\nano\farad}.
  \item \textbf{Nunca} conectar \SI{24}{\volt} directo al BNC, ni usar phantom de \SI{48}{\volt}~\cite{pcbguide}.
  \item \textbf{El amplificador} se alimenta con los \SI{3.3}{\volt} del microcontrolador. Las tierras (malla BNC, negativo de \SI{24}{\volt}, GND del micro)
        se unen en \textbf{un solo punto} (tierra en estrella) junto a los BNC.
  \item \textbf{El ruido dominante} a frecuencias de audio no es la fuente de corriente (su aporte, estimado en
        $\approx$\SI{4}{\micro\volt} rms, es menor que el ruido propio del micrófono), sino el \textbf{zumbido de \SI{60}{\hertz}}
        por lazos de tierra: osciloscopio a tierra física + fuente + USB del PC. Medidas: fuente flotante o baterías,
        portátil con batería al usar el USB, cable coaxial corto.
\end{itemize}

%=====================================================================
\section{Construcción sin PCB}

\textbf{¿Protoboard o no?} La capacitancia entre filas de una protoboard es de $\approx$\SI{2}{\pico\farad}~\cite{breadboard}; a
\SI{20}{\kilo\hertz} eso es $\approx$\SI{4}{\mega\ohm} de reactancia, despreciable frente a la salida de $<\SI{50}{\ohm}$ del micrófono.
Analog Devices desaconseja la protoboard para circuitos analógicos de alto desempeño~\cite{mt100}, pero para audio con fuente
de baja impedancia suele ser utilizable si se cuidan las tierras. Plan propuesto:
\begin{enumerate}
  \item \textbf{Prototipo en protoboard} con cables cortos, condensadores de desacople en cada op-amp, la protoboard
        dentro de una lata metálica conectada a tierra y los BNC fijos a la lata (no colgando de la protoboard). Se mide el ruido
        con la entrada en corto.
  \item \textbf{Si hay demasiado ruido}, se rehace sobre \textbf{placa de cobre virgen} en técnica \emph{Manhattan}~\cite{manhattan}:
        toda la placa es plano de tierra, los nodos se hacen con pequeños cuadros de cobre pegados, no hay que perforar, y son
        $\approx$20 soldaduras. Es más rápido y limpio que la perforada punto a punto.
  \item Para el op-amp: zócalo DIP pegado sobre la placa, o el op-amp invertido (\emph{dead bug}).
\end{enumerate}

%=====================================================================
\section{Plan de pruebas}

\begin{enumerate}
  \item \textbf{Fuente de corriente con carga ficticia.} Amperímetro en serie con una resistencia de \SI{3.3}{\kilo\ohm} ($\approx$ un
        micrófono a \SI{12}{\volt}). Variar la resistencia de \SI{0}{} a \SI{4.7}{\kilo\ohm} y verificar que $I_0$ no cambia más
        de unos pocos por ciento. Medir la tensión en circuito abierto ($\approx$\SI{22}{\volt}).
  \item \textbf{RG-50.} Conectar, esperar \SI{1}{\minute}, medir $V_B$ con multímetro (esperado \SI{12.5}{\volt}). Observar en el
        osciloscopio (AC) la respuesta a la voz y a un tono de \SI{1}{\kilo\hertz} de un celular.
  \item \textbf{SV 12 recuperado.} (a) Verificar si tiene cápsula: sin cápsula no responde al sonido; no tocar el pin central.
        (b) Conectar con adaptador TNC hembra--BNC macho a la fuente de \SI{2.5}{\milli\ampere}. (c) Medir $V_B$ tras \SI{1}{\minute}:
        un valor estable entre \SI{2}{} y \SI{20}{\volt} indica que está vivo (registrar el valor, no es público);
        $\approx$\SI{22}{\volt} indica circuito abierto/dañado; $\approx$\SI{0}{\volt} indica corto. (d) Prueba acústica con
        el celular. (e) Si no tiene cápsula: inyectar una senoidal de $\leq$\SI{1}{\volt} a través de \SI{18}{\pico\farad}
        (el ``dummy'' ST~02 de Svantek) y verificar la salida (atenuación $\approx$\SI{0.7}{\decibel})~\cite{svan977w}.
  \item \textbf{Cadena de ADC.} Con el generador de funciones (\SI{1}{\kilo\hertz}, \SI{50}{\milli\volt}) en lugar del
        micrófono, verificar $G=11$, nivel medio de \SI{1.65}{\volt} y la ausencia de recorte.
  \item \textbf{Ruido.} Entrada en corto: medir el piso de ruido y el zumbido de \SI{60}{\hertz} en el osciloscopio y en la FFT.
\end{enumerate}

%=====================================================================
\section{Lista de materiales y costo estimado}

Precios consultados en línea el 4 de noviembre de 2026; \textbf{verificar existencias} antes de ir
(Electrónicos Caldas: 321~896~5106). Los artículos sin precio se estiman.

\begin{table}[H]
\centering
\small
\begin{tabularx}{\textwidth}{Yllr}
\toprule
Artículo & Cant. & Proveedor & Costo [COP] \\
\midrule
TL431 & 2 & Electrónicos Caldas & 1.100 \\
Transistor PNP 2N3906 / BC557 & 2 (+2 repuesto) & Local (\nota{verificar}) & $\approx$2.000 \\
Condensador poliéster \SI{1}{\micro\farad} & 2 & Electrónicos Caldas & 900 \\
Resistencias, 1N4148, 1N4007, condensadores varios & --- & Electrónicos Caldas & $\approx$5.000 \\
MCP6002 (DIP-8) \textit{o} LM358 & 1 & Didácticas \textit{o} Caldas & 3.142 \textit{o} 1.100 \\
Conector BNC hembra de chasis & 2 & Didácticas (CON-035A) & 4.056 \\
Adaptador TNC hembra -- BNC macho (para el SV 12) & 1 & Didácticas & 3.928 \\
Interruptor SPDT de palanca & 1 & Local & $\approx$1.500 \\
Placa de cobre virgen $10\times5$ cm & 1 & Electrónicos Caldas & 1.500 \\
Raspberry Pi Pico \textit{(o ESP32 DevKitC)} & 1 & Didácticas \textit{(o Caldas)} & 19.897 \textit{(27.300)} \\
Caja: lata metálica reutilizada & 1 & --- & 0 \\
\midrule
\multicolumn{3}{l}{\textbf{Total aproximado} (sin envío)} & \textbf{$\approx$ 43.000} \\
\bottomrule
\end{tabularx}
\end{table}
Elementos del laboratorio: fuente de \SI{24}{\volt}, osciloscopio, multímetro, generador de funciones, cables BNC y T BNC, protoboard.

%=====================================================================
\section{Cronograma}

\begin{table}[H]
\centering
\small
\begin{tabularx}{\textwidth}{lY}
\toprule
Fecha & Actividad \\
\midrule
4--7 nov & Presentar esta propuesta y obtener aprobación. Identificar físicamente los micrófonos. \\
7--10 nov & Compra de componentes (recogida en Electrónicos Caldas + pedido a Didácticas). \\
10--13 nov & Montaje de las fuentes de corriente, prueba con carga ficticia, prueba del SV 12. \textbf{Hito 1: forma de onda en el osciloscopio.} \\
13--18 nov & Etapa de acople, ganancia y antialias; captura con el microcontrolador; script de Python. \textbf{Hito 2.} \\
18--21 nov & Traslado a placa de cobre / caja metálica, medición de ruido, ajustes. \\
21--24 nov & Informe y presentación. \\
\textbf{25 nov} & \textbf{Entrega.} \\
\bottomrule
\end{tabularx}
\end{table}

%=====================================================================
\section{Riesgos y preguntas abiertas}

\begin{itemize}
  \item \textbf{El SV 12 recuperado puede estar dañado o no tener cápsula.} Mitigación: el sistema se demuestra con el RG-50; el
        SV 12 se caracteriza eléctricamente (polarización y prueba con \SI{18}{\pico\farad}).
  \item \textbf{Datos no publicados del SV 12:} tensión de polarización y corriente máxima. Se usa \SI{2.5}{\milli\ampere}, valor
        documentado por Svantek.
  \item \textbf{Versión del RG-50:} si es antigua, el rango es \SIrange{2}{6}{\milli\ampere}; los \SI{3.9}{\milli\ampere} propuestos sirven igual.
  \item \textbf{Disponibilidad y tiempos de envío} desde Medellín: hay alternativas locales para casi todo (LM358, ESP32).
  \item \textbf{Zumbido de \SI{60}{\hertz}:} plan de tierras y blindaje descrito en las secciones 9 y 10.
\end{itemize}

\subsection{Extensión opcional: nivel de presión sonora}
\label{sec:extension}
Con la sensibilidad $S$ conocida, el nivel se obtiene de las muestras como
\begin{equation}
  L_p = 20\log_{10}\!\left(\frac{v_{\text{rms}}/(G\,S)}{\SI{20}{\micro\pascal}}\right)\ \si{\decibel}\ \text{SPL},
\end{equation}
y la ponderación A puede aplicarse con un filtro digital (por ejemplo, el del proyecto \texttt{esp32-i2s-slm}~\cite{slm}).

%=====================================================================
\begin{thebibliography}{99}
\small
\bibitem{svan977w} Svantek, \emph{SVAN 977W — Appendix C: Technical specifications}, 2017. \url{https://svantek.com/wp-content/uploads/2021/05/svan_977w_appendix_c_2017-03-06.pdf}
\bibitem{svan958a} Svantek, \emph{SVAN 958A User Manual} v1.02, 2023. \url{https://svantek.com/wp-content/uploads/2023/01/SVAN958A-Man.EN_v.1.02_2023-03-03.pdf}
\bibitem{sv977d} Svantek, \emph{SV 977D User Manual} v1.00, 2022. \url{https://svantek.com/wp-content/uploads/2022/06/SV-977D_Man_EN_v.1.00_2022-11-23.pdf}
\bibitem{sv971a} Svantek, \emph{SV 971A User Manual} v1.03, 2024. \url{https://svantek.com/wp-content/uploads/2023/01/SV971A-Man.EN_v.1.03_2024-05-28.pdf}
\bibitem{rg50} ROGA Instruments, \emph{RG-50 Measurement Microphone — Datasheet}. \url{https://roga-instruments.com/downloads/PDF/RG-50-EN.pdf}
\bibitem{roga14} ROGA Instruments, \emph{1/4'' measurement microphones}. \url{https://roga-instruments.com/1-4-inch-mic/}
\bibitem{pcbguide} PCB Piezotronics, \emph{General Signal Conditioning Guide} (G0001E). \url{https://www.pcb.com/techsupport/docs/pcb/PCB-G0001E-0303.pdf}
\bibitem{pcbcable} PCB Piezotronics / IMI, \emph{Tip from Tech: Driving Long Cable Lengths}. \url{https://www.pcb.com/contentstore/mktgcontent/IMI_Downloads/TipfromTech-DrivingLongCableLengths.pdf}
\bibitem{hbk} HBK, \emph{TN119: IEPE Sensing}, 2021. \url{https://www.hbm.com/fileadmin/mediapool/files/technical-articles-technotes-white-papers/2021/TN119_HBK_IEPE-Sensing.pdf}
\bibitem{tida010249} Texas Instruments, \emph{TIDA-010249: Four-channel IEPE vibration sensor interface} (TIDUF13). \url{https://www.ti.com/lit/pdf/TIDUF13}
\bibitem{tida01471} Texas Instruments, \emph{TIDA-01471: IEPE vibration sensor interface reference design for PLC analog input} (TIDUD62). \url{https://www.ti.com/lit/pdf/TIDUD62}
\bibitem{slyt768} Texas Instruments, \emph{High-side current sources for industrial applications}, Analog Design Journal (SLYT768). \url{https://www.ti.com/lit/pdf/SLYT768}
\bibitem{tl431} Texas Instruments, \emph{TL431 / TL432 Precision Programmable Reference — Datasheet}. \url{https://www.ti.com/lit/ds/symlink/tl431.pdf}
\bibitem{lm334} Texas Instruments, \emph{LM134/LM234/LM334 3-Terminal Adjustable Current Sources} (SNVS746E). \url{https://www.ti.com/lit/ds/symlink/lm334.pdf}
\bibitem{lm317} onsemi, \emph{LM317 Three-Terminal Adjustable Positive Voltage Regulator}. \url{https://www.onsemi.com/download/data-sheet/pdf/lm317-d.pdf}
\bibitem{cn0540} Analog Devices, \emph{CN0540: IEPE-compatible data acquisition for piezo sensors}. \url{https://www.analog.com/CN0540}
\bibitem{cn0582} Analog Devices, \emph{CN0582: IEPE sensor interface}. \url{https://www.analog.com/CN0582}
\bibitem{mt100} W.~Kester, J.~Bryant, \emph{MT-100: Breadboarding and Prototyping Techniques}, Analog Devices. \url{https://www.analog.com/media/en/training-seminars/tutorials/MT-100.pdf}
\bibitem{endevco} Endevco, \emph{Integrating Isotron (IEPE) accelerometers with your own current source}. \url{https://endevco.com/our-resources/ask-the-experts/integrating-isotron-iepe-accelerometers-with-your-own-current-source}
\bibitem{esp134} R.~Elliott, \emph{Project 134: 4 mA current loop microphone supply}, ESP. \url{https://www.sound-au.com/project134.htm}
\bibitem{breadboard} Wikipedia, \emph{Breadboard}. \url{https://en.wikipedia.org/wiki/Breadboard}
\bibitem{manhattan} \emph{Manhattan style construction}. \url{https://pic.hallikainen.org/techref/pcb/manhattan.htm}
\bibitem{picoex} Raspberry Pi, \emph{pico-examples: adc/dma\_capture}. \url{https://github.com/raspberrypi/pico-examples/tree/master/adc}
\bibitem{rp2040err} kitanokitsune, \emph{rp2040adc\_correction} (corrección del error E11 del ADC del RP2040). \url{https://github.com/kitanokitsune/rp2040adc_correction}
\bibitem{espadc} Espressif, \emph{ESP-IDF: ADC continuous mode driver}. \url{https://docs.espressif.com/projects/esp-idf/en/stable/esp32/api-reference/peripherals/adc/adc_continuous.html}
\bibitem{serialplot} H.~Y.~Özderya, \emph{SerialPlot}. \url{https://github.com/hyOzd/serialplot}
\bibitem{scoppy} \emph{Scoppy: oscilloscope and logic analyzer for the Raspberry Pi Pico}. \url{https://github.com/fhdm-dev/scoppy}
\bibitem{slm} I.~Kostoski, \emph{esp32-i2s-slm: sound level meter}. \url{https://github.com/ikostoski/esp32-i2s-slm}
\bibitem{caldas} Electrónicos Caldas, catálogo en línea. \url{https://www.electronicoscaldas.com/es/}
\bibitem{didacticas} Didácticas Electrónicas, catálogo en línea. \url{https://didacticaselectronicas.com/}
\end{thebibliography}

\end{document}
